\begin{exercise}{49}
Let $\mathcal C$ be a simplicial category. Show that the coherent nerve $\mathrm{N}(\mathcal C)$ is isomorphic
to the nerve of an ordinary category if and only if $\mathcal C$ is in the image of the functor $c\colon\mathrm{Cat}\to
\mathrm{Cat}_{\Delta}$.

\begin{proof}
One direction is Lemma~1.2.64. For the converse, suppose
$\mathrm{N}(\mathcal C)$ is isomorphic to an ordinary nerve, so its simplices
are uniquely determined by their spines (Theorem~1.1.52).
For an edge $\gamma\colon f\to g$ in a hom-simplicial set
$\mathrm{Hom}_{\mathcal C}(x,y)$, the triangle constructed in
Exercise~\ref{ex:42} has successive edges $\mathrm{id}_x,g$.
The constant hom-edge $s_0g$ gives another triangle with
the same spine. Uniqueness makes these triangles equal,
and hence $\gamma=s_0g$.

Suppose some hom-simplicial set $K=\mathrm{Hom}_{\mathcal C}(x,y)$
has a nondegenerate simplex of positive dimension.
Choose one $\eta\in K_m$ of least positive dimension;
then $m\ge2$.
By minimality and Exercise~\ref{ex:5}, every proper face is
totally degenerate, and these constants have a common vertex
$w$, since all edges of $\eta$ are degenerate.
Thus the map $\widehat\eta\colon\Delta^m\to K$ is constant
at $w$ on $\partial\Delta^m$.
Put $n=m+1$ and define
\[
 \chi\colon\mathrm{N}(P_{0,n})
 \xrightarrow{\ \mathrm{N}(\rho)\ }\Delta^m
 \xrightarrow{\ \widehat\eta\ }K,
 \qquad \rho(S)=|S|-2.
\]
Here we use Definition~1.2.56 and Lemma~1.1.44 to identify
$\mathrm{N}([m])$ with $\Delta^m$.
On any chain of subsets all containing a fixed
$k\in\{1,\ldots,m\}$, the rank map omits $0$;
therefore $\chi$ is constant at $w$ on that subcomplex.
But on the flag
\[
 \{0,n\}\subset\{0,1,n\}\subset\cdots\subset\{0,1,\ldots,n\}
\]
its value is $\eta$.

Construct two simplicial functors
$F,G\colon\mathfrak C[\Delta^n]\to\mathcal C$.
Both send $0,\ldots,m$ to $x$ and $n$ to $y$.
Their hom-maps between the first $m+1$ objects are constant
at $\mathrm{id}_x$, and those from these objects to $n$
are constant at $w$, except that $F_{0,n}=\chi$.
At $(n,n)$ both use the unit at $y$.
These assignments respect composition: every composition
through an intermediate $0<k<n$ into the exceptional hom
is induced by union of subsets containing $k$, where
$\chi$ is constant at $w$.
All other composition identities follow from the unit laws
and composition of constant simplices; empty homs impose no condition.
Thus $F,G$ are $n$-simplices of $\mathrm{N}(\mathcal C)$
with the same spine. They must be equal, whereas evaluating
their $(0,n)$ components on the displayed flag gives
$\eta$ and the totally degenerate $m$-simplex at $w$,
a contradiction.

Consequently every hom-simplicial set is constant on its
vertices, by Exercise~\ref{ex:5}.
Sending a vertex to its totally degenerate simplices gives
an isomorphism $c(u\mathcal C)\cong\mathcal C$:
it respects units and composition because the latter are
maps of simplicial sets. This proves the converse.
\end{proof}
\end{exercise}
